%=============================================================================
\section{Related Work}
\label{sec:related}
%=============================================================================

\paragraph{Spectral analysis and random matrix theory.}
The connection between eigenvalue spectra and neural network quality originates with the Marchenko-Pastur law~\citep{marchenko_pastur1967}, which characterizes random matrices and provides a null model for untrained weights.
Pennington \& Worah~\citep{pennington2017} extended free probability theory to nonlinear networks, showing that spectra deviate systematically from MP as learning progresses.
Martin \& Mahoney~\citep{martin_mahoney_jmlr2021} made the decisive discovery that well-trained layers exhibit heavy-tailed eigenvalue distributions whose power-law exponent $\alpha$ correlates with model quality, codified in the WeightWatcher tool~\citep{weightwatcher}.
Hodgkinson et al.~\citep{hodgkinson2025} recently provided theoretical underpinning by proposing random matrix ensembles that mechanistically explain why heavy tails emerge in weights, Jacobians, and Hessians of trained networks.
Complementary spectral perspectives include stable rank as a regularizer~\citep{sanyal2019stable}, the spectral edge thesis as a phase-transition framework~\citep{spectral_edge_thesis}, high-entropy advantages for generalization~\citep{high_entropy_npj2026}, spin-glass models connecting spectral structure to hidden order~\citep{barney_spin_glass2024, winter_janssen2025}, and scaling collapse revealing universal spectral dynamics across scales~\citep{scaling_collapse_icml2025}.
Our work extends all of these from \emph{static} post-hoc analysis to \emph{dynamic} real-time monitoring, introduces SR/$d$ as a complementary metric with stronger cross-scale correlation ($r = -0.92$) than $\alpha$ alone, and discovers the $\alpha$ reversal phenomenon and three-phase training dynamics.

\paragraph{Thermodynamics and statistical physics of training.}
A growing body of work frames neural network training through the lens of statistical physics.
Sadrtdinov et al.~\citep{sadrtdinov2025} showed that SGD implicitly minimizes a free energy functional; Liu et al.~\citep{liu_tegmark2025} identified learning rate with effective temperature via fluctuation-dissipation relations and derived thermodynamic laws for LLM training; the same group verified the ideal-gas equation $PV\!=\!NkT$ for scale-invariant networks at small scale~\citep{pvnkt_ijcai2026}.
Liu \& Chuang~\citep{liu_ziyin_neurips2025} developed a neural thermodynamics framework based on entropic forces, while Okanohara~\citep{okanohara2026} formalized irreversible learning as ensemble transport.
Phase transitions have been identified in grokking~\citep{grokking_glass_neurips2025, grokking_pretraining2025}, continual learning~\citep{order_params_cl2024}, LLM scaling~\citep{phase_transitions_on_model}, and diffusion model training~\citep{biroli_mezard2025_diffusion}.
Spring-block models connect layer-wise feature learning to mechanical relaxation~\citep{spring_block_prl2025}, and renormalization group methods reveal universal learning dynamics~\citep{rg_for_dnns2025}.
Our contribution is not a new state equation, but the identification of \emph{spectral} thermodynamic quantities (SR as $\exp(H_2)$, $\alpha$ as order parameter) that provide robust real-time diagnostics where macroscopic state equations fail at scale.

\paragraph{Scaling laws, training practice, and adaptive schedules.}
Kaplan et al.~\citep{kaplan2020} established power-law scaling of loss with model size and data, and Hoffmann et al.~\citep{hoffmann2022} refined this to compute-optimal ratios ($D/N \approx 20$).
Subsequent work demonstrated that training well beyond Chinchilla ratios yields superior downstream performance~\citep{gadre2024, llama3, muennighoff2023}, though overtrained models can be harder to fine-tune~\citep{overtraining_fragile2025}.
On the scheduling side, Wen et al.~\citep{wsd_schedule} analyzed Warmup-Stable-Decay through a river-valley loss landscape lens; theoretical foundations for optimal schedules have been developed under functional scaling laws~\citep{li_schedules2026, bordelon_schedules2026}; and annealing strategies have been studied for scaling and transferability~\citep{annealing_scaling_law2024, annealing_scaling_aaai2026}.
Current training practice monitors loss, gradient norms, and loss-landscape curvature~\citep{deepseekv3, kalra_curvature2026}, occasionally complemented by hyperparameter transfer via $\mu$P~\citep{mup} or learning dynamics analysis during continual pretraining~\citep{cpt_dynamics2025}.
Our Structural Chinchilla law provides a spectral explanation for why over-training works: structural maturity ($\alpha < 3$) requires $D/N > 500$, and the maturation timescale $\tau$ grows with model size.
Unlike fixed or hand-tuned schedules, our $\alpha$-guided approach derives the decay point from the model's own spectral state.
